% Gebrauchsinformation (patient information leaflet) for a FICTIONAL medicine "Beispielin 400 mg Filmtabletten".
% Test document, not a real leaflet: the product, the active substance and the company do not exist. German, 3 pages.
% Features: numbered sections 1-6 as in real leaflets, dosage table by age and body weight, side effects by frequency,
% warnings, missed dose, storage, contents of the pack.
\documentclass[9.5pt,a4paper]{extarticle}
\usepackage[T1]{fontenc}
\usepackage[utf8]{inputenc}
\usepackage[scaled=0.95]{helvet}
\renewcommand{\familydefault}{\sfdefault}
\usepackage[ngerman]{babel}
\usepackage[margin=2.0cm,top=2.3cm,bottom=2.3cm]{geometry}
\usepackage{xcolor}
\usepackage{tcolorbox}
\usepackage{enumitem}
\usepackage{booktabs}
\usepackage{array}
\usepackage{tabularx}
\usepackage{titlesec}
\usepackage{fancyhdr}
\usepackage{lastpage}
\usepackage{needspace}
\usepackage{microtype}

\definecolor{pharma}{HTML}{0B6E4F}
\setlength{\parindent}{0pt}
\setlength{\parskip}{0.45em}
\renewcommand{\arraystretch}{1.18}
\setlist{leftmargin=1.3em, itemsep=0.1em, topsep=0.2em}

\titleformat{\section}{\large\bfseries\color{pharma}}{\thesection.}{0.5em}{}[{\color{pharma}\titlerule}]
\titleformat{\subsection}{\bfseries}{}{0pt}{}
\titlespacing*{\section}{0pt}{1.1em}{0.4em}
\titlespacing*{\subsection}{0pt}{0.8em}{0.1em}
\renewcommand{\thesection}{\arabic{section}}

\pagestyle{fancy}
\fancyhf{}
\renewcommand{\headrulewidth}{0.4pt}
\fancyhead[L]{\small\textbf{Beispielin 400 mg Filmtabletten}}
\fancyhead[R]{\small Gebrauchsinformation: Information für Anwender}
\fancyfoot[L]{\scriptsize\textbf{TESTDOKUMENT:} fiktives Arzneimittel, keine echte Packungsbeilage, nicht zur Anwendung bestimmt.}
\fancyfoot[R]{\scriptsize Seite \thepage{} von \pageref{LastPage}}

\newcommand{\wichtig}[1]{\begin{tcolorbox}[colback=pharma!8, colframe=pharma, boxrule=0.6pt, arc=2pt, left=7pt, right=7pt,
  top=4pt, bottom=4pt]#1\end{tcolorbox}}

\begin{document}

% ---- Kopf ----------------------------------------------------------------------------------
{\Large\bfseries Gebrauchsinformation: Information für Anwender}\\[0.3em]
{\fontsize{20}{24}\selectfont\bfseries\color{pharma} Beispielin 400 mg Filmtabletten}\\[0.2em]
{\large Wirkstoff: Probanol}

\vspace{0.2em}
\wichtig{\textbf{Lesen Sie die gesamte Packungsbeilage sorgfältig durch, bevor Sie mit der Einnahme dieses
Arzneimittels beginnen, denn sie enthält wichtige Informationen.}
\begin{itemize}
  \item Heben Sie die Packungsbeilage auf. Vielleicht möchten Sie diese später nochmals lesen.
  \item Wenn Sie weitere Fragen haben, wenden Sie sich an Ihren Arzt oder Apotheker.
  \item Dieses Arzneimittel wurde Ihnen persönlich verschrieben. Geben Sie es nicht an Dritte weiter.
  \item Wenn Sie sich nach 3 Tagen nicht besser oder gar schlechter fühlen, wenden Sie sich an Ihren Arzt.
\end{itemize}}

\subsection*{Was in dieser Packungsbeilage steht}
\begin{enumerate}[label=\arabic*., itemsep=0pt]
  \item Was ist Beispielin und wofür wird es angewendet?
  \item Was sollten Sie vor der Einnahme von Beispielin beachten?
  \item Wie ist Beispielin einzunehmen?
  \item Welche Nebenwirkungen sind möglich?
  \item Wie ist Beispielin aufzubewahren?
  \item Inhalt der Packung und weitere Informationen
\end{enumerate}

% ---- 1 ------------------------------------------------------------------------------------------
\section{Was ist Beispielin und wofür wird es angewendet?}
Beispielin enthält den Wirkstoff Probanol. Probanol gehört zu einer Gruppe von Arzneimitteln, die Schmerzen lindern
und Fieber senken.

Beispielin wird angewendet zur kurzzeitigen Behandlung von leichten bis mäßig starken Schmerzen, wie
Kopfschmerzen, Zahnschmerzen oder Regelschmerzen, und von Fieber.

% ---- 2 ------------------------------------------------------------------------------------------
\section{Was sollten Sie vor der Einnahme von Beispielin beachten?}
\subsection*{Beispielin darf nicht eingenommen werden,}
\begin{itemize}
  \item wenn Sie allergisch gegen Probanol oder einen der in Abschnitt 6 genannten sonstigen Bestandteile dieses
        Arzneimittels sind;
  \item wenn Sie in den letzten 3 Monaten der Schwangerschaft sind (ab der 28.\ Schwangerschaftswoche);
  \item wenn Sie schon einmal ein Magen- oder Darmgeschwür oder eine Magen-Darm-Blutung hatten;
  \item wenn Sie an einer schweren Leber-, Nieren- oder Herzschwäche leiden;
  \item bei Kindern unter 12 Jahren oder mit einem Körpergewicht unter 40 kg.
\end{itemize}

\subsection*{Warnhinweise und Vorsichtsmaßnahmen}
Sprechen Sie vor der Einnahme von Beispielin mit Ihrem Arzt oder Apotheker, wenn Sie Asthma haben, Bluthochdruck
haben oder wenn Sie zu Magenbeschwerden neigen. Nebenwirkungen können verringert werden, indem die niedrigste
wirksame Dosis über den kürzesten für die Behandlung erforderlichen Zeitraum angewendet wird.

\wichtig{\textbf{Achtung:} Nehmen Sie Beispielin nicht länger als 3 Tage bei Fieber und nicht länger als 4 Tage bei
Schmerzen ein, ohne ärztlichen Rat einzuholen. Überschreiten Sie nicht die im Abschnitt 3 angegebene Höchstdosis.}

\subsection*{Einnahme von Beispielin zusammen mit anderen Arzneimitteln}
Informieren Sie Ihren Arzt oder Apotheker, wenn Sie andere Arzneimittel einnehmen oder kürzlich eingenommen
haben, insbesondere Arzneimittel zur Hemmung der Blutgerinnung, Arzneimittel gegen Bluthochdruck, entwässernde
Arzneimittel sowie weitere Schmerzmittel. Nehmen Sie nicht zwei Arzneimittel mit schmerzlindernder Wirkung
gleichzeitig ein, ohne Ihren Arzt zu fragen.

\subsection*{Schwangerschaft und Stillzeit}
Wenn Sie schwanger sind oder stillen, oder wenn Sie vermuten, schwanger zu sein, fragen Sie vor der Einnahme dieses
Arzneimittels Ihren Arzt oder Apotheker um Rat. In den ersten 6 Monaten der Schwangerschaft sollten Sie Beispielin
nur nach Rücksprache mit Ihrem Arzt einnehmen. In den letzten 3 Monaten der Schwangerschaft dürfen Sie
Beispielin nicht einnehmen.

\subsection*{Verkehrstüchtigkeit und Fähigkeit zum Bedienen von Maschinen}
Bei kurzzeitiger Anwendung hat Beispielin keinen oder nur einen vernachlässigbaren Einfluss auf die Verkehrstüchtigkeit.
Wenn bei Ihnen Schwindel auftritt, fahren Sie nicht Auto und bedienen Sie keine Maschinen.

\subsection*{Beispielin enthält Lactose}
Bitte nehmen Sie Beispielin erst nach Rücksprache mit Ihrem Arzt ein, wenn Ihnen bekannt ist, dass Sie unter einer
Unverträglichkeit gegenüber bestimmten Zuckern leiden.

% ---- 3 ------------------------------------------------------------------------------------------
\Needspace{20\baselineskip}
\section{Wie ist Beispielin einzunehmen?}
Nehmen Sie dieses Arzneimittel immer genau so ein, wie es in dieser Packungsbeilage beschrieben ist bzw. genau nach
Absprache mit Ihrem Arzt oder Apotheker. Fragen Sie bei Ihrem Arzt oder Apotheker nach, wenn Sie sich nicht sicher
sind. Die empfohlene Dosis richtet sich nach Alter und Körpergewicht:

\begin{center}
\small
\begin{tabularx}{\linewidth}{@{}>{\raggedright\arraybackslash}p{3.5cm} >{\raggedright\arraybackslash}p{2.0cm} >{\raggedright\arraybackslash}X >{\raggedright\arraybackslash}p{2.7cm} >{\raggedright\arraybackslash}p{3.0cm}@{}}
  \toprule
  \textbf{Alter} & \textbf{Körpergewicht} & \textbf{Einzeldosis} & \textbf{Abstand zwischen den Einnahmen} & \textbf{Höchstdosis in 24 Stunden} \\
  \midrule
  Kinder unter 12 Jahren & unter 40 kg & nicht empfohlen & --- & --- \\
  \addlinespace[3pt]
  Jugendliche von 12 bis 14 Jahren & 40 bis 49 kg & ½ Filmtablette (200 mg) & mindestens 8 Stunden & 3 halbe Filmtabletten (600 mg) \\
  \addlinespace[3pt]
  Erwachsene und Jugendliche ab 15 Jahren & ab 50 kg & 1 Filmtablette (400 mg) & mindestens 6 Stunden & 3 Filmtabletten (1\,200 mg) \\
  \bottomrule
\end{tabularx}
\end{center}

Die Filmtabletten werden unzerkaut mit reichlich Flüssigkeit, zum Beispiel einem Glas Wasser, eingenommen, am besten
während oder nach einer Mahlzeit. Die Tablette kann an der Bruchkerbe geteilt werden.

\subsection*{Dauer der Anwendung}
Beispielin ist nur zur kurzzeitigen Anwendung bestimmt. Nehmen Sie es bei Fieber nicht länger als 3 Tage und bei
Schmerzen nicht länger als 4 Tage ein, ohne mit Ihrem Arzt gesprochen zu haben.

\subsection*{Wenn Sie eine größere Menge von Beispielin eingenommen haben, als Sie sollten}
Wenn Sie zu viele Filmtabletten eingenommen haben, sprechen Sie sofort mit einem Arzt oder Apotheker oder wenden
Sie sich an die nächste Notaufnahme oder an eine Giftnotrufzentrale. Nehmen Sie diese Packungsbeilage und die
Packung mit.

\subsection*{Wenn Sie die Einnahme von Beispielin vergessen haben}
Nehmen Sie nicht die doppelte Menge ein, wenn Sie die vorherige Einnahme vergessen haben. Wenn Sie noch
Beschwerden haben, nehmen Sie die nächste Dosis ein, sobald Sie daran denken. Halten Sie dabei den in der
Tabelle angegebenen Abstand zur letzten Einnahme ein (mindestens 6 Stunden, bei Jugendlichen von 12 bis 14
Jahren mindestens 8 Stunden) und überschreiten Sie die Höchstdosis in 24 Stunden nicht. Wenn Sie keine Beschwerden
mehr haben, brauchen Sie die vergessene Dosis nicht nachzuholen.

% ---- 4 ------------------------------------------------------------------------------------------
\Needspace{24\baselineskip}
\section{Welche Nebenwirkungen sind möglich?}
Wie alle Arzneimittel kann auch dieses Arzneimittel Nebenwirkungen haben, die aber nicht bei jedem auftreten müssen.

\wichtig{\textbf{Beenden Sie die Einnahme und suchen Sie sofort einen Arzt auf,} wenn Sie eines der folgenden
Anzeichen bemerken: Schwellung von Gesicht, Zunge oder Rachen, Atemnot, Blut im Erbrochenen, schwarzer Stuhl oder
starke Bauchschmerzen. Dies können Zeichen einer schweren allergischen Reaktion oder einer Blutung im Magen-Darm-Trakt
sein.}

Die Häufigkeitsangaben sind wie folgt definiert:

\begin{center}
\small
\begin{tabularx}{\linewidth}{@{}>{\raggedright\arraybackslash}p{2.2cm} >{\raggedright\arraybackslash}p{4.4cm} >{\raggedright\arraybackslash}X@{}}
  \toprule
  \textbf{Häufigkeit} & \textbf{Das bedeutet} & \textbf{Beobachtete Nebenwirkungen} \\
  \midrule
  Sehr häufig & kann mehr als 1 von 10 Behandelten betreffen & Magenbeschwerden wie Sodbrennen, Übelkeit und Bauchschmerzen \\
  \addlinespace[3pt]
  Häufig & kann bis zu 1 von 10 Behandelten betreffen & Kopfschmerzen, Schwindel, Durchfall, Verstopfung, Blähungen, Erbrechen \\
  \addlinespace[3pt]
  Gelegentlich & kann bis zu 1 von 100 Behandelten betreffen & Hautausschlag, Juckreiz, Schlaflosigkeit, Unruhe, Ohrgeräusche, Entzündung der Mundschleimhaut \\
  \addlinespace[3pt]
  Selten & kann bis zu 1 von 1\,000 Behandelten betreffen & Magen-Darm-Geschwüre und -Blutungen, erhöhte Leberwerte, Wassereinlagerungen (Ödeme), schwere allergische Reaktionen mit Schwellung von Gesicht, Zunge oder Kehlkopf \\
  \bottomrule
\end{tabularx}
\end{center}

\subsection*{Meldung von Nebenwirkungen}
Wenn Sie Nebenwirkungen bemerken, wenden Sie sich an Ihren Arzt oder Apotheker. Dies gilt auch für Nebenwirkungen,
die nicht in dieser Packungsbeilage angegeben sind. Sie können Nebenwirkungen auch direkt dem Bundesinstitut für
Arzneimittel und Medizinprodukte melden (Internet: www.bfarm.de). Indem Sie Nebenwirkungen melden, können Sie dazu
beitragen, dass mehr Informationen über die Sicherheit dieses Arzneimittels zur Verfügung gestellt werden.

% ---- 5 ------------------------------------------------------------------------------------------
\section{Wie ist Beispielin aufzubewahren?}
Bewahren Sie dieses Arzneimittel für Kinder unzugänglich auf. Sie dürfen das Arzneimittel nach dem auf dem
Umkarton und dem Blister angegebenen Verfallsdatum (Verwendbar bis) nicht mehr verwenden. Das Verfallsdatum bezieht
sich auf den letzten Tag des angegebenen Monats. Nicht über 25\,\textdegree C lagern. In der Originalverpackung
aufbewahren, um den Inhalt vor Feuchtigkeit zu schützen. Entsorgen Sie Arzneimittel nicht im Abwasser oder
Haushaltsabfall. Fragen Sie Ihren Apotheker, wie das Arzneimittel zu entsorgen ist, wenn Sie es nicht mehr
verwenden.

% ---- 6 ------------------------------------------------------------------------------------------
\section{Inhalt der Packung und weitere Informationen}
\subsection*{Was Beispielin enthält}
Der Wirkstoff ist Probanol. Jede Filmtablette enthält 400 mg Probanol. Die sonstigen Bestandteile sind:
mikrokristalline Cellulose, Maisstärke, Lactose-Monohydrat, Croscarmellose-Natrium, hochdisperses Siliciumdioxid,
Magnesiumstearat (Ph.\,Eur.); Filmüberzug: Hypromellose, Titandioxid, Macrogol.

\subsection*{Wie Beispielin aussieht und Inhalt der Packung}
Beispielin sind weiße, ovale Filmtabletten mit einer Bruchkerbe auf einer Seite und der Prägung „BSP 400“. Es ist in
Packungen mit 10, 20 und 50 Filmtabletten erhältlich. Es werden möglicherweise nicht alle Packungsgrößen in den
Verkehr gebracht.

\subsection*{Pharmazeutischer Unternehmer und Hersteller}
Musterpharma GmbH, Beispielweg 12, 12345 Musterstadt (fiktive Firma)

\subsection*{Stand}
Diese Packungsbeilage wurde zuletzt überarbeitet im Oktober 2026.

\end{document}
